% !TeX root = ../tfg.tex
% !TeX encoding = utf8
%
% ====================================================================
%  CAPÍTULO 2 — TEMA 2: ANÁLISIS Y DISEÑO DE PUESTOS DE TRABAJO
%
%  Este capítulo cubre todos los apartados del Tema 2 de la asignatura:
%    2.1  Análisis de los puestos de trabajo (APT): etapas en CaixaBank
%    2.2  Descripción y especificación de tres puestos
%    2.3  Formas de rediseño de puestos
%    2.4  Nuevas tendencias: ADP por competencias
% ====================================================================

\chapter{Tema 2: Análisis y Diseño de Puestos de Trabajo}


% ================================================================
\section{Análisis de los puestos de trabajo en CaixaBank}
% ================================================================

% RECORDATORIO AL EQUIPO:
% Las etapas del análisis de puestos (APT) son generalmente:
%   1. Planificación del APT (objetivos, recursos)
%   2. Preparación (diseño de instrumentos: cuestionarios, entrevistas...)
%   3. Recogida de información (observación, entrevista, cuestionario, diario...)
%   4. Análisis y tratamiento de datos
%   5. Elaboración de la descripción y especificación del puesto
%   6. Actualización y mantenimiento
%
% PARA CADA ETAPA INDICAD:
%   - ¿La realiza CaixaBank? ¿Cómo exactamente?
%   - Ventajas e inconvenientes de su enfoque
%   - ¿Está la empresa satisfecha con ese proceso?
%   - ¿Habría maneras más eficientes de hacerlo?

El Análisis del Puesto de Trabajo (APT) es el proceso a través del cual la empresa recopila y analiza la información sobre el puesto de trabajo con la intención de identificar las tareas, obligaciones y sus responsabilidades, de forma que sirvan para establecer el perfil de persona que debería ocuparlo.

A continuación se detalla en qué medida CaixaBank desarrolla cada una
de las etapas clásicas del APT:

% ---- TABLA RESUMEN DE ETAPAS ----
\begin{table}[H]
  \centering
  \small
  \caption{Resumen de etapas del APT en CaixaBank}
  \begin{tabularx}{\textwidth}{p{0.30\textwidth} p{0.12\textwidth} X}
    \toprule
    \textbf{Etapa} & \textbf{¿Se realiza?} & \textbf{Observaciones} \\
    \midrule
    1. Planificación del APT          & Sí & (Liderado por Secretaría Técnica de Negocios y Deloitte). \\
    2. Preparación de instrumentos    & Sí & (Enfoque bottom-up, basado en las necesidades del cliente). \\
    3. Recogida de información        & Sí & (Adaptada a la ubicación y contexto de la oficina). \\
    4. Análisis y tratamiento         & Sí & (Apoyado por consultoría externa). \\
    5. Descripción y especificación   & Sí & (Totalmente orientado a competencias y soft skills). \\
    6. Actualización y mantenimiento  & Sí & (Dinámico y ajustado a la demanda del mercado). \\
    \bottomrule
  \end{tabularx}
\end{table}

% ---- ETAPA 1 ----
\subsection{Etapa 1 — Planificación del análisis de puestos}

% TODO: ¿Planifica formalmente CaixaBank el APT? ¿Tiene un departamento
% o equipo dedicado? ¿Con qué periodicidad lo actualiza?
% ¿Lo hace internamente o externaliza (consultoría)?

A diferencia de los modelos teóricos tradicionales, en CaixaBank la planificación del análisis de puestos no recae de forma exclusiva en el departamento de Recursos Humanos. Esta etapa está externalizada y liderada por la Secretaría Técnica de Negocios, utilizando un modelo desarrollado por la consultora externa Deloitte.

\textbf{Ventajas del enfoque adoptado:}
\begin{itemize}
  \item Alineación con el negocio: Al ser diseñado por la Secretaría Técnica, el puesto nace con un enfoque 100\% comercial y estratégico, garantizando que el empleado aporte valor real a la cuenta de resultados de la entidad.
  \item Objetividad y profesionalidad: Contar con una consultora experta (Deloitte) aporta metodologías probadas en el mercado y evita sesgos internos.
\end{itemize}

\textbf{Inconvenientes / áreas de mejora:}
\begin{itemize}
  \item Coste económico: Externalizar este proceso a consultoras de primer nivel supone un alto coste para la organización.
  \item Distanciamiento de RRHH: Que RRHH no defina el puesto puede generar desajustes a la hora de realizar la posterior selección o formación si la comunicación con la Secretaría Técnica no es fluida.
\end{itemize}

% ---- ETAPA 2 ----
\subsection{Etapa 2 — Preparación de los instrumentos de recogida}

% TODO: ¿Qué métodos de recogida de información utiliza CaixaBank?
% Métodos habituales: cuestionario estructurado, entrevista al ocupante,
% entrevista al supervisor, observación directa, diario de actividades,
% incidentes críticos, conferencia técnica de expertos.
% Indicad cuáles SÍ usa, cuáles NO y por qué (coste, tipo de puesto...).

CaixaBank utiliza un enfoque radicalmente distinto al tradicional: lo hacen "de abajo arriba" (bottom-up). En lugar de que la dirección imponga unas tareas cerradas, la recogida de información se realiza partiendo de las necesidades reales del cliente en la base.

Dado que el diseño lo elabora Negocio junto a Deloitte, se deduce que el método principal para recopilar datos no es la simple observación, sino los grupos de expertos. Este método es muy útil cuando lo que se pretende es el diseño de un nuevo puesto o un rediseño sustancial, reuniendo a directores de zona o especialistas para definir qué se necesita exactamente en la red comercial.

% ---- ETAPA 3 ----
\subsection{Etapa 3 — Recogida de información}

% TODO: ¿Quién recoge la información? ¿El propio empleado, su supervisor,
% un analista de RRHH? ¿Hay sesgo? ¿Cómo garantizan la fiabilidad?

Además, la recogida de información tiene muy en cuenta el contexto local, ya que el diseño del puesto y las habilidades requeridas varían "dependiendo del lugar donde está la oficina".

% ---- ETAPAS 4, 5 y 6 ----
\subsection{Etapas 4, 5 y 6 — Análisis, descripción y actualización}

% TODO: ¿Existe un catálogo formal de puestos? ¿Con qué frecuencia se revisa?
% En un banco como CaixaBank la aparición de nuevas figuras (analista de datos,
% ciberseguridad, UX bancario) exige actualizaciones frecuentes.
% ¿Lo hace CaixaBank? ¿Cómo gestionan la obsolescencia de los perfiles?

En la fase de descripción y especificación del puesto, CaixaBank aplica directamente las nuevas tendencias en RRHH: el análisis de puestos basado en competencias.

El modelo de Deloitte no redacta descripciones de puestos basadas únicamente en un listado rígido de tareas repetitivas, sino que implica la redacción de la descripción teniendo en cuenta los conocimientos, experiencias y, sobre todo, las competencias que se requieren para que se pueda desempeñar el puesto con éxito.

Le dan una importancia capital a las soft skills (competencias blandas como la empatía, la capacidad de negociación, la adaptabilidad o el trato al cliente). Esto permite que los puestos no se queden obsoletos rápidamente, ya que si cambian las herramientas informáticas del banco, el empleado con las soft skills adecuadas (flexibilidad, aprendizaje) se adaptará sin necesidad de tener que rehacer por completo la descripción formal de su puesto.


% ================================================================
\section{Descripción y especificación de tres puestos de CaixaBank}
% ================================================================

% RECORDATORIO AL EQUIPO:
% Hay que presentar TRES puestos distintos de la empresa.
% Para cada puesto debéis elaborar:
%
% A) DESCRIPCIÓN DEL PUESTO (el QUÉ se hace):
%    - Identificación: denominación, departamento, nivel jerárquico,
%      superior inmediato, subordinados, fecha de elaboración.
%    - Resumen del puesto (misión en 2-3 líneas).
%    - Funciones y tareas principales (ordenadas por importancia/frecuencia).
%    - Responsabilidades: personas, equipos, presupuesto, información.
%    - Condiciones de trabajo: horario, lugar, movilidad, riesgos.
%    - Relaciones internas y externas.
%
% B) ESPECIFICACIÓN DEL PUESTO (el QUIÉN lo puede hacer):
%    - Formación mínima requerida.
%    - Experiencia mínima.
%    - Conocimientos específicos.
%    - Competencias / habilidades (técnicas y conductuales).
%    - Rasgos de personalidad relevantes.
%
% PARA CADA PUESTO VALORAD:
%    - ¿La descripción está actualizada? ¿Refleja la realidad?
%    - ¿La especificación es adecuada (ni sobre- ni infracualificación)?
%    - ¿Está la empresa satisfecha con estas fichas?
%    - ¿Podría mejorarse? ¿Cómo?
%
% Elegid tres puestos de diferentes niveles jerárquicos o áreas:
%   - Director/a de Oficina (nivel directivo / mando intermedio)
%   - Gestor/a de Banca Privada (nivel técnico especializado)
%   - Asesor/a Comercial de Oficina (nivel operativo)


% ============================================================
\subsection{Puesto 1: Director/a de Oficina}
% ============================================================

\subsubsection{Descripción del puesto}

\begin{table}[H]
  \centering
  \small
  \caption{Ficha de identificación — Puesto 1}
  \begin{tabularx}{\textwidth}{lX}
    \toprule
    \textbf{Campo}                & \textbf{Contenido} \\
    \midrule
    Denominación del puesto       & Director/a de Oficina \\
    Código / referencia interna   & [Si existe] \\
    Departamento / Área           & Red Territorial / Oficinas \\
    Nivel jerárquico              & Mando intermedio \\
    Superior inmediato            & Director/a Territorial de Zona \\
    Número de subordinados        & Entre 3 y 10 personas según el tamaño de la Store u oficina \\
  \end{tabularx}
\end{table}

\textbf{Misión del puesto:}

Gestionar, coordinar y liderar el equipo de la oficina bancaria, asegurando el cumplimiento de los objetivos comerciales, la rentabilidad de la sucursal y la excelencia en la satisfacción del cliente, operando siempre bajo los estrictos estándares de riesgo y cumplimiento normativo del Grupo CaixaBank.

\textbf{Funciones y tareas principales:}
\begin{enumerate}
  \item Planificar y asignar los objetivos comerciales individuales al equipo de la oficina.
  \item Realizar el seguimiento sistemático de la actividad comercial y proponer planes de acción ante desviaciones.
  \item Gestionar directamente la cartera de clientes de mayor valor o riesgo de la sucursal.
  \item Liderar el desarrollo profesional, la motivación y la evaluación del desempeño de los empleados a su cargo.
  \item Garantizar el cumplimiento de las normativas de auditoría interna, prevención de blanqueo de capitales y MiFID II.
\end{enumerate}

\textbf{Responsabilidades:}
\begin{itemize}
  \item \textit{Sobre personas:} Liderazgo, supervisión, evaluación y motivación de un equipo (habitualmente entre 3 y 10 personas según el tamaño de la Store u oficina).
  \item \textit{Sobre equipos / materiales:} Gestión de la infraestructura de la oficina y seguridad física.
  \item \textit{Sobre información confidencial:} Acceso a datos patrimoniales de clientes, gestión de riesgos crediticios y morosidad.
  \item \textit{Sobre presupuesto:} Responsabilidad directa sobre la cuenta de resultados (P\&L) de la oficina y el cumplimiento del presupuesto anual asignado por la Dirección Territorial.
\end{itemize}

\textbf{Condiciones de trabajo:}
\begin{itemize}
  \item Horario: Jornada completa, habitualmente horario partido (mañana y tarde, especialmente en el modelo "Oficina Store" de CaixaBank).
  \item Lugar: Oficina física, con escasa posibilidad de teletrabajo.
  \item Movilidad / desplazamientos: Sí, ocasionales para visitas a empresas clientes o reuniones en la Dirección Territorial.
  \item Riesgos laborales: Alto nivel de estrés por presión comercial y cumplimiento de objetivos; riesgos psicosociales y ergonomía por uso de pantallas.
\end{itemize}

\subsubsection{Especificación del puesto}

\begin{table}[H]
  \centering
  \small
  \caption{Requisitos del candidato — Puesto 1}
  \begin{tabularx}{\textwidth}{lX}
    \toprule
    \textbf{Requisito}        & \textbf{Descripción} \\
    \midrule
    Formación mínima          & Grado en ADE, Economía, Derecho o similar \\
    Formación complementaria  & Certificación MiFID II (imprescindible) y Ley de Contrato de Crédito Inmobiliario (LCCI) \\
    Experiencia mínima        & Al menos 5 años en el sector bancario, con un mínimo de 2-3 años de experiencia previa en gestión de equipos o como Subdirector/a \\
    Conocimientos técnicos    & Análisis de riesgos de empresas y particulares, cuenta de resultados, productos de inversión y financiación \\
    Idiomas                   & Español nativo; nivel alto del idioma cooficial (si aplica a la CCAA); inglés nivel B2 valorable \\
    Competencias conductuales & Liderazgo transformacional, toma de decisiones bajo presión, alta orientación a resultados y resiliencia \\
    Herramientas / software   & Terminal Financiero propio de CaixaBank, CRM interno, Office 365 \\
    \bottomrule
  \end{tabularx}
\end{table}

En general, la especificación interna del puesto de Director/a de Oficina es más exigente que las ofertas externas en:
\begin{itemize}
  \item \textbf{Reclutamiento:} las ofertas externas tienden a centrarse en habilidades comerciales y cumplimiento de KPI; la ficha interna refuerza el enfoque en riesgo y cumplimiento.
  \item \textbf{Requisitos formativos:} el grado y la certificación MiFID II son estándar, pero la inclusión de LCCI y experiencia en gestión de equipos puede reducir el pool de candidatos.
  \item \textbf{Competencias:} la ficha interna prioriza resiliencia y liderazgo; las ofertas externas suelen insistir más en orientación comercial y resultados.
\end{itemize}

% ============================================================
\subsection{Puesto 2: Gestor/a de Banca Privada}
% ============================================================

\subsubsection{Descripción del puesto}

\begin{table}[H]
  \centering
  \small
  \caption{Ficha de identificación — Puesto 2}
  \begin{tabularx}{\textwidth}{lX}
    \toprule
    \textbf{Campo}                & \textbf{Contenido} \\
    \midrule
    Denominación del puesto       & Gestor/a de Banca Privada \\
    Departamento / Área           & Banca Privada / CaixaBank Private Banking \\
    Nivel jerárquico              & Técnico especializado \\
    Superior inmediato            & Director/a de Banca Privada \\
    Número de subordinados        & Ninguno / gestión de cartera propia \\
    \bottomrule
  \end{tabularx}
\end{table}

\textbf{Misión del puesto:}

Proporcionar un asesoramiento patrimonial y financiero integral y altamente personalizado a clientes del segmento de rentas altas (Banca Privada), buscando la optimización fiscal y la rentabilidad de sus inversiones para fidelizar e incrementar los activos bajo gestión (AUM).

\textbf{Funciones y tareas principales:}
\begin{enumerate}
  \item Diseñar, presentar y ejecutar propuestas de inversión personalizadas acordes al perfil de riesgo del cliente.
  \item Realizar un seguimiento proactivo de los mercados financieros y rebalancear las carteras de los clientes cuando sea necesario.
  \item Captar nuevo patrimonio y nuevos clientes dentro del segmento objetivo.
  \item Coordinar con especialistas internos (fiscalistas, gestores de fondos) la estructuración del patrimonio del cliente.
\end{enumerate}

\textbf{Responsabilidades:}
\begin{itemize}
  \item \textit{Sobre personas:} Ninguna. Es un perfil de contribuidor individual.
  \item \textit{Sobre equipos / materiales:} Equipos informáticos corporativos en movilidad (portátil, smartphone).
  \item \textit{Sobre información confidencial:} Confidencialidad extrema sobre grandes patrimonios, estructuras societarias y datos sensibles.
  \item \textit{Sobre presupuesto:} Responsabilidad sobre el crecimiento del volumen de negocio (Assets Under Management) de su cartera asignada.
\end{itemize}

\textbf{Condiciones de trabajo:}
\begin{itemize}
  \item Horario: Jornada completa con alta flexibilidad horaria para adaptarse a las necesidades de clientes VIP.
  \item Lugar: Centro de Banca Privada y alta movilidad.
  \item Movilidad / desplazamientos: Muy frecuente (visitas a domicilios, empresas de clientes o eventos de networking).
  \item Riesgos laborales: Estrés ligado a la volatilidad de los mercados financieros.
\end{itemize}

\subsubsection{Especificación del puesto}

\begin{table}[H]
  \centering
  \small
  \caption{Requisitos del candidato — Puesto 2}
  \begin{tabularx}{\textwidth}{lX}
    \toprule
    \textbf{Requisito}        & \textbf{Descripción} \\
    \midrule
    Formación mínima          & Grado en ADE, Economía o Ciencias Actuariales \\
    Certificaciones           & Certificación de Asesor Financiero (Nivel EFPA o equivalente) obligatoria; CFA muy valorable \\
    Experiencia mínima        & Entre 3 y 5 años de experiencia demostrable en gestión de patrimonios o mercados financieros \\
    Conocimientos técnicos    & Profundo conocimiento macroeconómico, fiscalidad de inversiones, SICAVs, fondos de inversión alternativos \\
    Competencias conductuales & Excelentes habilidades de comunicación, empatía, discreción absoluta, proactividad y tolerancia a la frustración \\
    \bottomrule
  \end{tabularx}
\end{table}


En comparación con el puesto de Director/a de Oficina, el Gestor/a de Banca Privada muestra:
\begin{itemize}
  \item \textbf{Mayor especialización técnica:} trabaja con productos complejos (SICAVs, fondos alternativos, fiscalidad patrimonial) y requiere certificaciones como EFPA/CFA.
  \item \textbf{Relación cliente más exclusiva:} gestiona carteras de clientes VIP y ofrece soluciones a medida, mientras que el director combina atención cliente con gestión de oficina.
  \item \textbf{Responsabilidad individual:} no supervisa equipo ni P\&L de la oficina, pero tiene responsabilidad directa sobre el crecimiento de los activos bajo gestión.
\end{itemize}

\textbf{Conclusión:} CaixaBank distingue claramente entre banca personal (oficina) y banca privada, manteniendo roles con distintos niveles técnicos y tipos de cliente, aunque ambos mantienen el foco en la venta y en la gestión del riesgo.


% ============================================================
\subsection{Puesto 3: Asesor/a Comercial de Oficina}
% ============================================================

\subsubsection{Descripción del puesto}

\begin{table}[H]
  \centering
  \small
  \caption{Ficha de identificación — Puesto 3}
  \begin{tabularx}{\textwidth}{lX}
    \toprule
    \textbf{Campo}                & \textbf{Contenido} \\
    \midrule
    Denominación del puesto       & Asesor/a Comercial de Oficina \\
    Departamento / Área           & Red Comercial \\
    Nivel jerárquico              & Operativo / técnico base \\
    Superior inmediato            & Director/a de Oficina \\
    Número de subordinados        & Ninguno \\
    Fecha de elaboración / revisión & [Fecha] \\
    \bottomrule
  \end{tabularx}
\end{table}

\textbf{Misión del puesto:}

Atender de forma proactiva y excelente a los clientes en la sucursal, resolviendo sus operativas diarias y detectando oportunidades para comercializar productos y servicios financieros y aseguradores básicos, derivando al cliente a perfiles especializados cuando sea necesario.

\textbf{Funciones y tareas principales:}
\begin{enumerate}
  \item Recepción, atención y derivación de clientes en la sucursal (labor de "filtro" o bienvenida).
  \item Venta cruzada de productos transaccionales, tarjetas, seguros de auto/hogar y préstamos al consumo.
  \item Resolución de incidencias operativas y soporte en el uso de los canales digitales (cajeros, app CaixaBankNow).
  \item Realización de campañas comerciales telefónicas (telemarketing) a la cartera de clientes asignada.
\end{enumerate}

\textbf{Responsabilidades:}
\begin{itemize}
  \item \textit{Sobre personas:} Ninguna.
  \item \textit{Sobre equipos / materiales:} Terminal financiero.
  \item \textit{Sobre información confidencial:} Protección de datos personales (RGPD) de los clientes de la oficina.
  \item \textit{Sobre presupuesto:} Cumplimiento de objetivos comerciales individuales (número de tarjetas, pólizas de seguro, etc.).
\end{itemize}

\textbf{Condiciones de trabajo:}
\begin{itemize}
  \item Horario: Jornada continua (normalmente de 8:00 a 15:00) o partida en oficinas Store, dependiendo del contrato.
  \item Lugar: Puesto fijo en la zona de atención al público de la sucursal.
  \item Movilidad / desplazamientos: Nula o muy baja.
  \item Riesgos laborales: Fatiga vocal, exposición al público (riesgo de clientes conflictivos), sedentarismo.
\end{itemize}

\subsubsection{Especificación del puesto}

\begin{table}[H]
  \centering
  \small
  \caption{Requisitos del candidato — Puesto 3}
  \begin{tabularx}{\textwidth}{lX}
    \toprule
    \textbf{Requisito}        & \textbf{Descripción} \\
    \midrule
    Formación mínima          & Grado Universitario (preferiblemente en rama económica/empresarial) o FP de Grado Superior en Administración y Finanzas \\
    Certificaciones           & Certificación MiFID II (muchas veces el banco asume esta formación tras la contratación) \\
    Experiencia mínima        & 0 a 2 años (es un puesto habitual de entrada o junior) \\
    Conocimientos técnicos    & Conocimiento de productos financieros básicos, ofimática y agilidad en sistemas informáticos \\
    Competencias conductuales & Gran orientación al cliente, dotes comerciales, empatía, paciencia y capacidad de trabajo en equipo \\
    \bottomrule
  \end{tabularx}
\end{table}


  \textbf{Rotación y ajuste del perfil:}
  \begin{itemize}
    \item \textbf{Rotación alta en nivel junior:} suele existir una brecha entre las expectativas de desarrollo y la realidad operativa; se puede mejorar con formación y acompañamiento inicial más estructurado.
    \item \textbf{Dificultad para cubrir vacantes:} las ofertas pueden describir un puesto más amplio del que luego resulta ser en el día a día; esto alarga el proceso de selección y genera fricción.
    \item \textbf{Desajuste perfil-puesto:} si se prioriza actitud comercial en detrimento de competencias técnicas básicas, los nuevos empleados pueden frustrarse rápidamente.
  \end{itemize}

  \textbf{Sugerencia:} Fortalecer el onboarding con formación práctica, clarificar objetivos comerciales desde el inicio y revisar incentivos para mejorar la retención.



% ================================================================
\section{Formas de rediseño de puestos en CaixaBank}
% ================================================================

% RECORDATORIO AL EQUIPO:
% Las principales formas de rediseño de puestos son:
%
%   1. ROTACIÓN DE PUESTOS (job rotation): los empleados rotan entre puestos.
%      Ventajas: polivalencia, motivación, reducción de la monotonía.
%      Inconvenientes: costes de formación, pérdida de especialización.
%
%   2. AMPLIACIÓN DE TAREAS (job enlargement / ensanchamiento horizontal):
%      se añaden tareas del mismo nivel de complejidad.
%      Ventajas: reduce la monotonía, mayor variedad.
%      Inconvenientes: puede aumentar la carga sin añadir motivación real.
%
%   3. ENRIQUECIMIENTO DE PUESTOS (job enrichment / Hackman & Oldham):
%      se añaden responsabilidades de mayor complejidad (carga vertical).
%      Las cinco características centrales del modelo:
%        - Variedad de habilidades
%        - Identidad de la tarea
%        - Significado de la tarea
%        - Autonomía
%        - Retroalimentación
%      Ventajas: mayor motivación intrínseca, satisfacción, compromiso.
%      Inconvenientes: requiere empleados que quieran / puedan asumir más.
%
%   4. EQUIPOS DE TRABAJO AUTOGESTIONADOS (self-managed teams).
%   5. TELETRABAJO / TRABAJO FLEXIBLE (flexibilización del puesto).
%   6. JOB CRAFTING: el empleado adapta el puesto a sus fortalezas
%      de forma proactiva (tendencia emergente).
%
% PARA CADA FORMA INDICAD:
%   - ¿La aplica CaixaBank? ¿En qué puestos / áreas?
%   - Ventajas e inconvenientes concretos observados.
%   - ¿Está la empresa satisfecha?
%   - ¿Qué formas serían interesantes aplicar y por qué?

\subsection{Formas de rediseño aplicadas actualmente}

\begin{table}[H]
  \centering
  \small
  \caption{Aplicación de formas de rediseño de puestos en CaixaBank}
  \begin{tabularx}{\textwidth}{p{0.30\textwidth} p{0.12\textwidth} X}
    \toprule
    \textbf{Forma de rediseño}       & \textbf{¿Se aplica?} & \textbf{Descripción / observaciones} \\
    \midrule
    Rotación de puestos              & [Sí/No/Parcial] & \\
    Ampliación de tareas             & [Sí/No/Parcial] & \\
    Enriquecimiento de puestos       & [Sí/No/Parcial] & \\
    Equipos autogestionados          & [Sí/No/Parcial] & \\
    Teletrabajo / trabajo flexible   & [Sí/No/Parcial] & \\
    Job crafting                     & [Sí/No/Parcial] & \\
    \bottomrule
  \end{tabularx}
\end{table}

\subsubsection{Rotación de puestos (\textit{job rotation})}

% TODO: ¿Existe un programa formal de rotación en CaixaBank?
% Es habitual en banca para desarrollar directivos y gestores con visión
% global. ¿Se rota entre oficinas, entre departamentos, internacionalmente
% (BPI)?

[Descripción a completar]

\textbf{Ventajas observadas en CaixaBank:}
\begin{itemize}
  \item [...]
\end{itemize}

\textbf{Inconvenientes observados:}
\begin{itemize}
  \item [...]
\end{itemize}

\subsubsection{Ampliación de tareas (\textit{job enlargement})}

% TODO: Con la digitalización bancaria, los asesores de oficina han
% asumido tareas que antes hacían otros departamentos (apertura online,
% soporte digital...). ¿Esto genera sobrecarga? ¿Hay compensación
% retributiva?

[Descripción a completar]

\subsubsection{Enriquecimiento de puestos (\textit{job enrichment})}

% TODO: ¿Se da autonomía a los asesores para tomar decisiones sobre
% productos? ¿Tienen feedback claro sobre su impacto en los resultados?
% El modelo de Hackman y Oldham es muy útil para evaluar si el diseño
% actual de los puestos genera motivación intrínseca.

[Descripción a completar]

\subsubsection{Teletrabajo y trabajo flexible}

% TODO: Tras la pandemia COVID-19, CaixaBank implementó políticas de
% teletrabajo. ¿Sigue vigente? ¿Para qué puestos? ¿Con qué porcentaje
% de jornada? ¿Ha mejorado la satisfacción y productividad?

[Descripción a completar]

\subsection{Formas de rediseño recomendadas para CaixaBank}

% TODO: Proponed al menos 2-3 formas que no se apliquen o que se apliquen
% de forma mejorable. Justificad por qué serían beneficiosas.

[Descripción a completar]

\begin{notaequipo}
  Preguntad en la entrevista:
  \begin{enumerate}
    \item ¿Existe algún programa formal de rotación? ¿Tiene criterios claros?
    \item ¿Los empleados de oficina han asumido más tareas en los últimos años?
          ¿Se ha revisado su carga de trabajo y retribución?
    \item ¿Tienen autonomía para adaptar su forma de trabajo (\textit{job crafting})?
    \item ¿Cuál es la política actual de teletrabajo? ¿Están satisfechos con ella?
    \item ¿Se mide el impacto de los rediseños en productividad y satisfacción?
  \end{enumerate}
\end{notaequipo}


% ================================================================
\section{Nuevas tendencias en análisis y diseño de puestos: el enfoque por competencias}
% ================================================================

% RECORDATORIO AL EQUIPO:
% El ADP por competencias supone un cambio de paradigma:
%
% ADP TRADICIONAL:
%   - Se centra en TAREAS y RESPONSABILIDADES del puesto.
%   - El resultado es una descripción de puesto estática.
%   - El encaje persona-puesto se hace comparando requisitos formales
%     (titulación, experiencia) con el perfil del candidato.
%
% ADP POR COMPETENCIAS:
%   - Se centra en los COMPORTAMIENTOS OBSERVABLES que llevan al éxito.
%   - Identifica competencias clave (técnicas + conductuales/transversales).
%   - Es más flexible: la misma competencia puede ser válida en varios puestos.
%   - Facilita la gestión integrada del talento bajo un marco común.
%   - Permite adaptar rápidamente los perfiles a entornos cambiantes (VUCA).
%
% OTRAS TENDENCIAS:
%   - ADP en entornos ágiles (puestos definidos por roles y sprints).
%   - Uso de People Analytics e IA para rediseñar puestos.
%   - Gestión por proyectos vs. gestión por puestos (modelos gig/fluid).
%   - ADP continuo: revisión periódica sistemática.
%   - Wellbeing y diseño de puestos orientado al bienestar.
%
% PARA ESTA SECCIÓN INDICAD:
%   - ¿Aplica CaixaBank el ADP por competencias?
%   - ¿Tiene un diccionario de competencias corporativo?
%   - ¿Cómo integra el enfoque en selección, formación y evaluación?
%   - ¿Usa People Analytics?

\subsection{Del enfoque tradicional al enfoque por competencias}

El análisis y diseño de puestos ha evolucionado desde el modelo clásico,
centrado en la descripción de tareas y responsabilidades, hacia un enfoque
basado en \textbf{competencias}, que pone el foco en los comportamientos
observables y medibles que conducen a un desempeño excelente.

\begin{table}[H]
  \centering
  \small
  \caption{Comparativa entre el ADP tradicional y el ADP por competencias}
  \begin{tabularx}{\textwidth}{p{0.23\textwidth} X X}
    \toprule
    \textbf{Dimensión}       & \textbf{ADP Tradicional} & \textbf{ADP por Competencias} \\
    \midrule
    Unidad de análisis       & Tareas y responsabilidades & Competencias y comportamientos \\
    Resultado                & Descripción de puesto estática & Perfil de competencias dinámico \\
    Orientación              & Presente                  & Presente y futuro \\
    Flexibilidad             & Baja                      & Alta \\
    Integración con RRHH     & Limitada                  & Total (marco integrado) \\
    Adaptación al cambio     & Reactiva                  & Proactiva \\
    \bottomrule
  \end{tabularx}
\end{table}

\subsection{El modelo de competencias de CaixaBank}

% TODO: CaixaBank, como gran empresa del IBEX 35, es probable que tenga
% un diccionario de competencias corporativo. Investigad:
%   - ¿Existe un modelo de competencias formalizado?
%   - ¿Cuáles son las competencias corporativas (transversales a todos)?
%   - ¿Hay competencias específicas por nivel o familia de puestos?
%   - ¿Cómo se usa el modelo en selección, evaluación y formación?

[Descripción a completar con la información de la entrevista y/o la
web corporativa de CaixaBank (sección de empleo / cultura)]

Ejemplos de competencias frecuentes en entidades bancarias
(ajustar a las reales de CaixaBank):
\begin{itemize}
  \item \textbf{Orientación al cliente}: escucha activa, empatía,
        resolución de problemas del cliente.
  \item \textbf{Orientación a resultados}: planificación, seguimiento
        de objetivos, toma de decisiones basada en datos.
  \item \textbf{Trabajo en equipo y colaboración}: comunicación, confianza,
        gestión de conflictos.
  \item \textbf{Adaptabilidad e innovación}: apertura al cambio,
        aprendizaje continuo, pensamiento creativo.
  \item \textbf{Integridad y cumplimiento}: ética profesional, respeto
        normativo, confidencialidad.
\end{itemize}

\subsubsection{Integración del modelo de competencias en los procesos de RRHH}

% TODO: ¿Cómo se usa el diccionario de competencias en:
%   - Selección (entrevista por competencias, assessment centres)?
%   - Evaluación del desempeño (evaluación 360º, por objetivos + competencias)?
%   - Formación y desarrollo (planes de desarrollo individuales)?
%   - Planes de carrera y sucesión?

[Descripción a completar]

\subsection{Otras tendencias emergentes}

\subsubsection{People Analytics aplicado al análisis de puestos}

% TODO: ¿Usa CaixaBank datos para analizar patrones de desempeño,
% rotación o satisfacción por tipo de puesto? ¿Tienen un equipo de
% People Analytics? ¿Qué sistemas usan (SAP, Workday, Cornerstone...)?

[Descripción a completar]

\subsubsection{ADP en entornos ágiles}

% TODO: Especialmente relevante en el área de tecnología y digitalización
% de CaixaBank. ¿Trabajan con metodologías ágiles (Scrum, SAFe)?
% ¿Cómo afecta eso al diseño de puestos (roles vs. puestos fijos)?

[Descripción a completar]

\subsubsection{Diseño de puestos y bienestar del empleado (\textit{wellbeing})}

% TODO: CaixaBank ha publicado informes de sostenibilidad con compromisos
% en bienestar del empleado. ¿Se refleja eso en el diseño de los puestos?
% ¿Tienen indicadores de bienestar vinculados al diseño de puestos?

[Descripción a completar]

\begin{notaequipo}
  Preguntad en la entrevista:
  \begin{enumerate}
    \item ¿Tienen un diccionario o modelo de competencias corporativo?
          ¿Podéis verlo o que os lo describan?
    \item ¿Cómo se usa ese modelo en la selección y en la evaluación
          del desempeño?
    \item ¿Tienen un equipo o función de People Analytics?
    \item ¿Trabajan con metodologías ágiles en alguna área? ¿Cómo ha
          cambiado el concepto de ``puesto'' en esas áreas?
    \item ¿Qué tendencias en ADP creen que les falta adoptar?
  \end{enumerate}
\end{notaequipo}
